\documentclass[10pt,a4paper]{amsart}
\usepackage[margin=19mm]{geometry}
\usepackage{amsmath,amssymb,mathtools}
\usepackage[scr=rsfs,cal=euler]{mathalfa}
\usepackage{xfrac}
\usepackage[unicode,hidelinks,bookmarks=false]{hyperref}
\newcommand{\ie}{i.e.\ }
\newcommand{\cf}{cf.\ }
\newcommand{\resp}{respectively\ }
\newcommand{\Subsection}{\subsection}
\providecommand{\nobreakdash}{\nobreak\hbox{-}}
\setlength{\emergencystretch}{2em}
\multlinegap0pt
\usepackage{amssymb}
\usepackage{tikz}
\newtheorem{theorem}{Theorem}
\title{Diophantine \emph{m}-tuples of Triangular Numbers}
\author{Sounak Bagchi, Christian Zhou-Zheng}
\date{Source: arXiv:2608.27697v1 (CC BY 4.0)}
\begin{document}
\maketitle

\noindent\emph{Source and licence.} Diophantine \emph{m}-tuples of Triangular Numbers. Version arXiv:2608.27697v1 (CC BY 4.0), distributed by arXiv under the Creative Commons Attribution 4.0 International licence, http://creativecommons.org/licenses/by/4.0/, as recorded in the arXiv licence metadata for this exact version and shown on the abstract page https://arxiv.org/abs/2608.27697v1. Full licence notice: \texttt{licenses/arxiv-2608.27697-CC-BY-4.0.txt}.\par\smallskip

\noindent\textit{Editorial scope.} Counted unit: the article's theorem theorem (source0.tex lines 436--438). The article's complete original proof of it is delivered with it (source0.tex lines 439--450, one \texttt{proof} environment of 12 lines). No mathematics is written, repaired or re-derived: the statement and the proof are the article's own lines, in the article's own notation, and no retained line is substituted or re-flowed.  No further source-local context is needed: every cross-reference of every retained line resolves inside the two delivered spans.  Nothing was omitted from any retained span, so the package carries no omission record.  No retained line contains a citation, so the wrapper carries no bibliography and no bibliography entry.  No retained line contains a figure environment, an image-inclusion command or drawing code, so no image is delivered and the package declares no asset dependency.  Editorial formatting only, in addition: an AMS \texttt{amsart} preamble that restates the article's own declarations and macros used by the retained lines, namely the environments \texttt{theorem} on separate counters, together with the standard packages those lines need.  The retained environments print in this document as Theorem 1 in order of appearance, whereas the article prints the same items with the numbering of its own full document.  The complete native source of the article as distributed in the arXiv e-print for this exact version is bundled unchanged as \texttt{sources/208731/source0.tex}.
\begin{theorem}
    The equation $T_a T_b + n = c^2$ has solutions $(a,b,c)$ in $\mathbb{Z} / 2^k \mathbb{Z}$, with $k \ge 1$, for all values of $n$.
\end{theorem}

\begin{proof}
    We prove a smaller claim first.
\\

    \textbf{Claim:} For $0 \le k \le 2^{m}-1$, $T_k$ modulo $2^m$ takes on every value in the set $\{0, 1, \dots, 2^m - 1\}$. 
    \\

    \emph{Proof:} Suppose that this is a contradiction, and that two values of $a,b$ exist for which $T_a \equiv T_b \pmod{2^m}$ and $a \neq b$. Then, it follows that $$\frac{a(a+1)}{2} - \frac{b(b+1)}{2} \equiv \frac{(a-b)(a+b+1)}{2} \equiv 0 \pmod{2^m}.$$ Note that $a-b$ and $a+b+1$ have opposite parity, so either $a \equiv b \pmod{2^{m+1}}$ or $a + b \equiv 2^{m+1} - 1 \pmod{2^{m+1}}$. The latter cannot be true since the maximum sum possible is $2 \left(2^{m} - 1 \right) = 2^{m+1} - 2$, hence $a \equiv b \pmod{2^m}$. So, $T_k$ has a unique residue modulo $2^m$ for $0 \le k \le 2^{m} - 1$, hence implying the conclusion. \quad\quad\quad\quad\quad\quad\quad\quad\quad\quad\quad\quad\quad\quad\quad\quad\quad $\square$
    \\ 
   
    Our claim makes the conclusion obvious, as $T_a T_b$ can take on any value. For example, set $T_a \equiv 4-n \pmod{2^k}$ and $T_b \equiv 1 \pmod{2^k}$, so that $T_a T_b + n \equiv (4-n) + n \equiv 2^2 \pmod{2^k}$. 
\end{proof}

\end{document}
